\documentclass[11pt,a4paper]{article}
\usepackage[utf8]{inputenc}
\usepackage[T1]{fontenc}
\usepackage{amsmath, amssymb, amsthm}
\usepackage{geometry}
\geometry{margin=2.5cm}
\usepackage{lmodern}
\usepackage[final,expansion=false]{microtype}
\usepackage{booktabs}
\usepackage{array}
\usepackage{enumitem}
\usepackage{xcolor}
\usepackage{titlesec}
\PassOptionsToPackage{hyphens}{url}
\usepackage{hyperref}
\usepackage{xurl}
\hypersetup{
  colorlinks=true,
  linkcolor=blue!50!black,
  urlcolor=blue!50!black,
  citecolor=blue!50!black,
  pdftitle={The Universal Right-Triangle Reflection Sequence -- Companion (conjectures and partial progress)},
  pdfauthor={Vico Bonfioli},
}
\usepackage{parskip}
\setlength{\parindent}{0pt}
\setlength{\parskip}{0.55em}

\titleformat{\section}{\normalfont\Large\bfseries\color{blue!40!black}}{\thesection.}{0.6em}{}

\newtheorem{conjecture}{Conjecture}
\newtheorem{proposition}{Proposition}
\theoremstyle{remark}
\newtheorem{remark}{Remark}
\newtheorem{observation}{Numerical observation}
\theoremstyle{plain}

\title{\textbf{Companion document: partial results, conjectures, and computational evidence}\\[0.3em]
\large Supplement to \emph{The Universal Right-Triangle Reflection Sequence}}
\author{Vico Bonfioli}
\date{\today}

\begin{document}
\maketitle

\paragraph{Preface.}
This companion document collects partial results, conjectures, and
computational evidence supporting the program of the main paper
\texttt{paper.tex}.  \textbf{Nothing here is claimed as a theorem.}
All results are explicitly conditional or empirical: each is labelled
either as a \emph{conjecture} (a target statement we believe but have
not proved), as a \emph{conditional proposition} (a deduction valid
under a hypothesis that is itself only empirically verified at a finite
horizon), or as a \emph{numerical observation} (a finite computational
record).  Cross-references to theorems of the main paper use the same
labels.

\section{The transcendence program for $\beta_2$: conjecture and three attack routes}
\label{sec:transcendence-program}

Define the asymptotic ratio
\[
\beta_2 \;:=\; \lim_{d \to \infty} \frac{u_{d+1}}{u_d}
\]
of the universal right-triangle reflection sequence.  The limit exists
and equals
\[
  \beta_2 \;=\; 1.4916177871143742268\ldots \;=\; 1/\sqrt{q^\ast},
\]
where $q^\ast = 0.44945363055894\ldots$ is the least positive root of
$\Sigma_1(q) = 1$; this is unconditional, and the singularity is a
simple pole ($\Sigma_0(q^\ast) \neq 0$).  Only the arithmetic nature of
$\beta_2$ remains open.

\begin{conjecture}[Transcendence of $\beta_2$]
\label{conj:beta2-transcendence}
$\beta_2$ is transcendental over $\mathbb{Q}$, and the generating
function $U(t) = \sum_{d \ge 0} u_d t^d$ is transcendental over
$\mathbb{Q}(t)$.
\end{conjecture}

We collect three independent attack routes.  Each rests on a different
finite-horizon hypothesis that is empirically verified and structurally
attackable; the routes are logically independent and would each
individually settle Conjecture~\ref{conj:beta2-transcendence}.

\subsection{Route 1: Christol mod-$p$ via the explicit cyclic-module structure}
\label{ssec:christol-route}

Let $\rho_{\mathrm{abs}}: W \to Q \ltimes M$ be the universal
realisation of the main paper (\S4) with $Q = D_\infty \times \mathbb{Z}/2$
and $M = \mathbb{Z}[Q]/(Y+1, c+1)$.  For each prime $p$, write
\[
S_p \;:=\; \bigl\{\,(\pi(g),\, m \bmod p) \;:\; (g,m) = \rho_{\mathrm{abs}}(w),\ w \in W \,\bigr\}
\;\subset\; Q \times (M/pM).
\]

\begin{proposition}[Mod-$p$ state space is infinite]
\label{prop:modp-state-infinite}
For every prime $p$, $|S_p| = \infty$.  In particular the infinite
family $\{((XY)^n, t_n) : n \ge 1\}$ (with $t_n$ the translation part
of $\rho_{\mathrm{abs}}((R_0 R_2)^n) \bmod p$) consists of pairwise
distinct elements of $Q \times (M/pM)$.
\end{proposition}

\emph{Sketch.}  $XY$ has infinite order in $D_\infty$, so the
$Q$-components $(XY)^n$ are pairwise distinct.  This alone suffices for
infinity of $S_p$ via projection.  The stronger separation on the
$(M/pM)$-coordinate, used in Route~1 below, follows from the witness
monomial $(XY)^{n-1} X \cdot h$ surviving reduction modulo $p$ by the
cyclic-module relations $(Y+1)h = (c+1)h = 0$.

\paragraph{Numerical record.}
An exact-arithmetic Rust BFS in $\mathbb{F}_p[Q]$ enumerates $S_p(d) :=
\{\rho_{\mathrm{abs}}(w) \bmod p : |w|_W \le d\}$ at three primes:
\[
\begin{array}{c|c|c|c}
p & d_{\max} & |S_p(d_{\max})| & \beta_p \text{ (ratio fit, last 10 depths)} \\ \hline
3 & 42 & 65\,428\,964 & 1.4681 \\
5 & 40 & 62\,736\,544 & 1.5048 \\
7 & 40 & 65\,724\,148 & 1.5077
\end{array}
\]
The implementation uses the fixed-degree truncation justified by
$(Y+1)h = (c+1)h = 0$, which bounds the support of any reachable
$m \in M/pM$ by $O(d)$ monomials.

\begin{proposition}[$U$ transcendence: Christol route, conditional]
\label{prop:beta2-christol}
Conditional on the Myhill--Nerode separation hypothesis (the states
associated to $\{((XY)^n, t_n)\}_{n \ge 1}$ in the $p$-DFAO for
$u_d \bmod p$ are pairwise distinct future-equivalence classes), $U(t)$
is transcendental over $\mathbb{Q}(t)$.
\end{proposition}

\emph{Sketch.}  Under the separation hypothesis the $p$-DFAO has
infinitely many states, hence (Allouche--Shallit Thm.~6.6.2) the
$p$-kernel of $U(t) \bmod p$ is infinite \cite{AlloucheShallit}, so by
Christol \cite{Christol1979,CKMR} $U(t) \bmod p$ is not algebraic over
$\mathbb{F}_p(t)$.  Adamczewski--Bell 2017 lifts this to non-algebraicity of $U(t)$ over $\mathbb{Q}(t)$.
The conclusion is transcendence of $U$, and it stops there;
see Remark~\ref{rem:fs-direction}.

\subsection{Route 2: D-finite exclusion}
\label{ssec:dfinite-route}

\begin{proposition}[No D-finite recurrence at low complexity]
\label{prop:no-dfinite}
Let $(k,m)$ satisfy $k \le 9$, $m \le 7$ and the over-determination
condition
\begin{equation}\label{eq:overdet}
   (k+1)(m+1) \;<\; 43 - k ,
\end{equation}
the right side being the number of equations available from
$u_0, \ldots, u_{42}$ at order $k$. There are $52$ such pairs, and on
none of them does a D-finite (P-recursive) recurrence of order $k$ with
polynomial-in-$d$ coefficients of degree $m$ exist (exact integer linear
algebra modulo the Mersenne prime $2^{61}-1$).
\end{proposition}

\begin{remark}[Why the guard is not optional]
\label{rem:overdet}
Condition~\eqref{eq:overdet} is what makes a negative result mean
anything. Of the $72$ pairs with $k \le 9$ and $m \le 7$, exactly $20$
violate it, and on those the linear system is underdetermined, so a
solution exists for \emph{any} input sequence whatever and no exclusion
can be read off. The extreme case is $(k,m) = (9,7)$, with $80$ unknowns
against $34$ equations. An earlier version of this proposition quantified
over all $47$ tested pairs without the guard and was therefore false as
stated; the companion~\cite{Bonfioli2} imposes the corresponding
condition explicitly. The same remark applies to any extension of the
search: raising $k$ or $m$ without also lengthening the sequence buys
nothing.

The underlying fact is elementary and is machine-checked in
\texttt{lean/with\_mathlib/OverDetermination.lean}. A homogeneous system
with strictly more unknowns than equations has a nontrivial solution for
\emph{every} coefficient matrix (\texttt{exists\_nonzero\_solution}),
so a search over such a shape reports the existence of a relation
whatever the input data are. Contrapositively, an exclusion that is
correct forces $\text{unknowns} \le \text{equations}$
(\texttt{le\_of\_exclusion}); that is the guard, and its failure convicts
the claim. The same file proves the two counts quoted here,
$\texttt{vacuous\_card} = 20$ of the $72$ pairs and
$252 > 38$ for the nonlinear shape, and records that the guard is
necessary but not sufficient (\texttt{guard\_not\_sufficient}): passing
it still leaves the rank of the coefficient matrix to be computed.
\end{remark}

\begin{proposition}[$U$ transcendence: D-finite route, conditional]
\label{prop:beta2-dfinite}
Conditional on extending Proposition~\ref{prop:no-dfinite} to all
$(k, m) \in \mathbb{N}^2$, $U(t)$ is transcendental over $\mathbb{Q}(t)$.
\end{proposition}

\begin{remark}[The direction of the Flajolet--Sedgewick step, and what these routes do not give]
\label{rem:fs-direction}
Earlier versions of both routes ended with the sentence that
Flajolet--Sedgewick, \emph{Analytic Combinatorics} Thm.~VII.7, ``implies
$\beta_2$ transcendental''.  That inference is backwards. Theorem~VII.7
says that an algebraic generating function has finitely many
singularities, all algebraic branch points or poles, so that
\emph{algebraic $U$ forces algebraic $\beta_2$}.  Contrapositively it
lets one deduce transcendence of $U$ from transcendence of $\beta_2$,
which is the opposite of what is wanted here.

The converse is false, so no repair of the wording will do. Take
\[
   U_\ast(t) \;=\; \frac{1}{1-2t} \;+\; \sum_{d \ge 0} t^{2^d}
   \;\in\; \mathbb{Z}[[t]] .
\]
Its coefficients are $2^n$ plus the indicator that $n$ is a power of
two, so the growth rate is $2$, an algebraic number, while the lacunary
part gives the unit circle as a natural boundary and $U_\ast$ is
transcendental over $\mathbb{Q}(t)$. Transcendence of the series is thus
compatible with an algebraic growth rate.

Consequently neither route above reaches $\beta_2$, and neither does
Route~3. What they reach, conditionally, is transcendence of $U$, and
that is now known unconditionally by the companion
paper~\cite{Bonfioli2}, which supersedes them. Transcendence of
$\beta_2$ itself remains open, and this note contains no approach to it.
\end{remark}

\emph{Sketch.}  By Stanley \emph{Enumerative Combinatorics} Vol.~2
Thm.~6.4.6, every power series algebraic over $\mathbb{Q}(t)$ is
D-finite.  Excluding D-finite at all $(k, m)$ therefore excludes
algebraic, and gives transcendence of $U$; it says nothing about
$\beta_2$, by Remark~\ref{rem:fs-direction}.  The natural
structural target for the extension is a Gr\"obner-basis argument on
the skew-polynomial ring $\mathbb{Q}\langle n, S\rangle$ acting against
the cyclic-module relations $(Y+1)h = (c+1)h = 0$.

\subsection{Route 3: Mahler structure}
\label{ssec:mahler-route}

\begin{conjecture}[$U(t)$ is not a class-(1) Mahler function]
\label{conj:not-class1-mahler}
$U(t)$ is not a Mahler function of class (1) in the sense of
Adamczewski--Faverjon 2020; for the general theory of Mahler functions
see \cite{Mahler1929,Nishioka}, and for the transcendence machinery that
a positive answer would feed, \cite{Nesterenko1996}.
\end{conjecture}

Were Conjecture~\ref{conj:not-class1-mahler} established, the Adamczewski--Faverjon
framework would yield a third independent route to transcendence of
$\beta_2$, parallel to Propositions~\ref{prop:beta2-christol} and~\ref{prop:beta2-dfinite}.

\subsection{Joint picture}

The three conditionals are logically independent, and all three bear on
the transcendence of $U$ rather than of $\beta_2$. Since
\cite{Bonfioli2} now settles transcendence of $U$ unconditionally, they
are retained only as a record of routes that were explored. The second
half of Conjecture~\ref{conj:beta2-transcendence}, transcendence of
$\beta_2$, is untouched by any of them.

\section{Structural exclusions for $U(t)$ at low complexity}
\label{sec:structural-negatives}

The following finite linear-algebra exclusions on $u_0, \ldots, u_{42}$
constrain the form of $U(t)$ at the verifiable horizon.

\begin{proposition}[No low-order linear recurrence]
\label{prop:no-recurrence}
The sequence $u_0, \ldots, u_{38}$ satisfies no linear recurrence over
$\mathbb{Q}$ of order $\le 19$.  (Equivalently: for every $k \le 19$,
the system $u_n = \sum_{j=1}^k c_j u_{n-j}$ has empty rational solution
set on the available terms.)
\end{proposition}

\begin{proposition}[No nonlinear or modular structure at low complexity]
\label{prop:no-recurrence-strong}
On $u_0, \ldots, u_{42}$:
\begin{enumerate}
\item[\textnormal{(i)}] No D-finite recurrence with $k \le 9$ and
$m \le 7$ (Proposition~\ref{prop:no-dfinite}).
\item[\textnormal{(ii)}] \emph{Withdrawn.} An earlier version asserted
that no nonlinear polynomial relation
$\sum_\alpha q_\alpha(n)\,u_n^{\alpha_0}\cdots u_{n-w}^{\alpha_w} = 0$
holds with lookback $w \le 5$, total $u$-degree $|\alpha| \le 3$ and
coefficient degree $\le 2$. That search is vacuous: the monomials of
total degree at most $3$ in six variables number
$\binom{9}{3} = 84$, so with three coefficient degrees there are $252$
unknowns against the $38$ equations available from
$u_0,\ldots,u_{42}$. An underdetermined system has solutions for every
input, so nothing is excluded. No nonlinear relation of that shape is
testable at all on $43$ terms.
\item[\textnormal{(iii)}] No non-trivial modular linear recurrence at
any modulus $m \in \{3, 4, \ldots, 101\}$ over the tested range, and at
$m = 2$ the order-$2$ Fibonacci recurrence
$u_d \equiv u_{d-1} + u_{d-2} \pmod 2$ holds for $d \ge 3$ but
\emph{not} at $d = 2$, where $u_0 + u_1 = 4 \equiv 0$ while
$u_2 = 5 \equiv 1$.
\end{enumerate}
\end{proposition}

\begin{proposition}[Coefficient-height gap]
\label{prop:height-gap}
$\log u_d = c \cdot d + o(d)$ with $c = \log \beta_2 = 0.39986129\ldots$,
this being forced by the growth rate established in \cite{Bonfioli2};
ordinary least squares on $d \in [25,42]$ returns $0.407927$, with
$R^2 \ge 0.999$, the gap from $\log\beta_2$ being the usual finite-horizon
bias. An earlier version quoted $c \approx 0.404$, which matches neither:
$e^{0.404} = 1.49780\ldots$ against $\beta_2 = 1.49161\ldots$.

The consequence to be drawn is only the weak one. Since $u_d$ grows
exponentially it is unbounded, so it is not the value sequence of any
$k$-automatic integer sequence, automatic sequences taking finitely many
values; no height-gap theorem is needed for that. In particular the
earlier appeal to a height gap to exclude Mahler functions was
misplaced: by Adamczewski, Bell and Smertnig \cite{ABS} the coefficient sequences of
Mahler functions realise five distinct growth behaviours, all of them
attained, so exponential growth excludes nothing.
\end{proposition}

\begin{proposition}[Mod-$2$ algebraicity]
\label{prop:mod2-automaton}
$u_d \bmod 2$ is \emph{eventually} $3$-periodic, with pre-period $1$:
$u_d \equiv 1$ for $d \not\equiv 0 \pmod 3$ or $d \le 2$, and
$u_d \equiv 0$ otherwise. It is not purely periodic, since $d = 0$ is
$\equiv 0 \pmod 3$ yet $u_0 \equiv 1$. The order-$2$ Fibonacci
recurrence $u_d \equiv u_{d-1} + u_{d-2} \pmod 2$ holds for $d \ge 3$,
and fails at $d = 2$: $u_0 + u_1 = 4 \equiv 0$ while $u_2 = 5 \equiv 1$.
Hence
$U(t) \in \mathbb{F}_2[[t]]$ is algebraic over $\mathbb{F}_2(t)$ via a
$4$-state finite automaton.

\emph{Caveat.}  This is a characteristic-$2$ statement and does not
lift to algebraicity over $\mathbb{Q}(t)$.
\end{proposition}

\section{Universality class: scope and known examples}
\label{sec:universality-class}

The mechanism that proves universality in the main paper
(the universality theorem of \cite{Bonfioli2}) (a cyclic-ideal identification
$I_A \cap I_B = \mathbb{Z}[Q]\cdot h$ with explicit central generator
$h$ inside the linear quotient $Q$ of a free product $W = A * B$) admits a natural class of free-product reflection systems.

\begin{conjecture}[Open question: scope of the cyclic-ideal mechanism]
\label{conj:universality-class}
Characterise the pairs $(W, \rho)$ for which the cyclic-ideal mechanism
of the universality theorem of \cite{Bonfioli2} applies and yields universality of
the orbit-growth sequence across the family of geometric realisations.
\end{conjecture}

Two examples are currently known.

\paragraph{Example A: right triangles.}  $A = V_4$, $B = \mathbb{Z}/2$,
$Q = D_\infty \times \mathbb{Z}/2$, $h = (c-1)(Y-1)$.  This is the main
paper's the universality theorem of \cite{Bonfioli2}.

\paragraph{Example B: hyperbolic two-right-angle quadrilaterals.}
$A = B = V_4$, $W = V_4 * V_4$, $Q = D_\infty \times \mathbb{Z}/2$,
$h = c - 1$ (no $(Y-1)$ factor needed: centrality of $c$ alone suffices
when both factors of the free product contain $c$). This example is
recorded as a claim only: we have neither a proof that the mechanism
applies nor a computational artifact for it, and it should not be relied
on.

\paragraph{Cases where the mechanism is structurally blocked.}  Acute
and obtuse triangle reflection groups have no central involution
playing the role of $c$.  Equilateral and other triangle groups with
all finite Coxeter labels lack the free-product structure.  The
rectangle group $D_\infty \times D_\infty$ is a direct product, not a
free product.  Right-corner orthoschemes in dimension $n \ge 3$ carry
cross-relations $(R_i R_j)^2 = 1$ between generators of distinct
free-product factors, again destroying the $W = A * B$ structure;
universality in this cyclic-ideal sense is therefore specific to
dimension~2.

\section{Structural reading of universality (algebraic-channel BFS)}
\label{sec:univ-algebraic}

Fix $\rho_{a,b}: W \to \mathrm{Aff}(\mathbb{R}^2)$.  Let $Y$ denote the
linear part of $\rho_{a,b}(R_2)$ and $v_2$ its translation part.  The
identity $Y v_2 = -v_2$ holds for every $(a, b)$ and is the image of a
single canonical word $w_\star \in W$.  Let $N_{\rm alg}$ be the normal
closure of $w_\star$ in $W$.

\begin{observation}[Empirical algebraic form of universality]
\label{obs:alg-structural}
Through depth $22$, on the four test right triangles $(3, 4), (5, 12),
(1, 2), (7, 24)$, and within the horizon in which the companion's
universality statement itself holds (it is proved to depth $32$ and
fails at depth $33$ for $(1,2)$), the equivalence relation on canonical Coxeter words
generated by $Y v_2 + v_2$ in $\mathbb{Z}[D_\infty]$ exactly recovers
the affine equivalence relation.  Equivalently: $N(T) = N_{\rm alg}$ on
the depth-$\le 22$ filtration for every tested $T$.
\end{observation}

A parallel BFS verifies this with cumulative state count $118\,938$ at
depth $22$ in every channel: four channels for the four test triangles
above and a fifth for the algebraic model $N_{\rm alg}$ against which
they are compared.  Reading: the eight
length-$10$ affine relations of the symbolic-verification theorem of \cite{Bonfioli2} are not
sporadic length-$10$ coincidences but the first eight consequences of
the single quadratic relation $Y v_2 + v_2 = 0$, made visible by the
canonical-word filtration.

\section{Empirical evidence against finite confluent rewriting}
\label{sec:rewriting}

\paragraph{Knuth--Bendix experiment.}
GAP \texttt{KBMAG} ran Knuth--Bendix completion on $W$ together with
the eight length-$10$ affine relations of the symbolic-verification theorem of \cite{Bonfioli2};
the procedure produced $32\,730$ derived rules in approximately $8$
seconds of \texttt{kbprog} computation before hitting the equation-count
cap, never reaching confluence.  This is empirical evidence (not a
proof) that $W/N$ has no finite confluent rewriting system in shortlex
on the standard generators.

\paragraph{Antol\'in--Ciobanu route (structurally blocked).}
Antol\'in and Ciobanu prove that the conjugacy-growth series of a
non-elementary acylindrically hyperbolic group is \emph{transcendental}
\cite{AntolinCiobanu}; an earlier version of this paragraph called it a
rationality theorem, which inverts the result, and attributed it to the
wrong journal. Either way it does not apply here, and for a reason that
is about the group and not about the series: the image of
$W/\ker(\rho_{\mathrm{abs}})$ lies inside
$Q \ltimes M \subseteq (\mathbb{Z}\wr\mathbb{Z})\rtimes V_4$, a virtually
metabelian group, which is not acylindrically hyperbolic. Note also that
$\mathbb{Z}\wr\mathbb{Z}$ is \emph{not} finitely presented, by Baumslag;
the earlier description of the ambient group as finitely presented was
wrong, and is not needed for the conclusion.

\section{Class A and Class B in 3D}
\label{sec:classes-AB}

The 3D right-corner orthoscheme reflection groups fall into three
faithfulness regimes by leg-symmetry. The Class~C case, pairwise distinct
legs, is treated in the $n$-dimensional companion~\cite{Bonfioli1b},
where the exclusion of rank-two relations is proved and faithfulness
itself is stated as a conjecture; it is not a theorem, and nothing below
assumes it.

\paragraph{Class A ($a = b = c$, $S_3$-symmetric, polynomial growth).}
The $(1, 1, 1)$ cube-corner orthoscheme is the fundamental chamber of
the affine Weyl group $\widetilde{C}_3$.  For $n \ge 1$,
$u_n = (8 n^2 + c_{n \bmod 5})/5$ with $c = (10, 12, 13, 13, 12)$, with
asymptotic $u_n \sim (8/5) n^2$.  The closed form, its integrality and
its agreement with the Bott series of $\widetilde{C}_3$ hold for
$n \ge 1$; in particular $u_1 = 4$ and $u_2 = 9$. Order-$7$ linear
recurrence with characteristic polynomial $(x-1)^3 \Phi_5(x)$, obtained
from the series $W_A(t) = N(t) / ((1-t)^3 \Phi_5(t))$ with
$N(t) = 2 + 3t^2 + 3t^3 + 3t^4 + t^5 + 4t^6$. Two cautions about that
display. Its constant term is $N(0) = 2$, whereas an orbit-growth series
has constant term $1$, so as written it is not the growth series but
that series plus one; the closed form for $n \ge 1$ is unaffected, and
the discrepancy sits entirely in degree $0$. And ``parabolic Poincar\'e
series'' is a misnomer for it. This is a clean
restatement of standard affine-Weyl-group material (Bourbaki Ch.~VI~§4;
Humphreys \S4.7) and is recorded here for completeness.

\paragraph{Class B (exactly two equal legs, intermediate exponential rate).}
Through depth $18$ the sequence on $(1, 1, 2)$ is
$1, 4, 9, 18, 35, 67, 129, 249, 480, 925, 1783, 3437, 6625, 12770,
24591, 47298, 90948, 174939, 336544$, with empirical ratio settling
around $\approx 1.92$.  No order-$\le 9$ linear recurrence over
$\mathbb{Q}$ fits the depth-$18$ data.  No closed form is known.

\section{Other deferred material}
\label{sec:deferred}

The following items were deferred from the main paper to this
supplement and remain available in earlier drafts of the project:
\begin{itemize}[leftmargin=*,itemsep=0.15em]
\item Numerical and structural analysis of the representation gap
$\beta_n$ vs.\ $r_n$ across dimensions.
\item Collision-class enumeration to depth $30$, the structural
sequence battery, the $\sqrt{3}$-asymptotic of $\delta_d - 2 z_d$, and
the crossover $\delta_d > u_d$ for $d \ge 24$.
\item $\{3, 4, 5\}$-cascade observation for maximal outerplanar graphs
on $n \le 15$ vertices, the slack/denominator/precision analysis, and
the cascade-existence conjecture.
\item The standard 3-simplex remark and merged length-$10$-relation
cascades from Pythagorean triples sharing a common edge length.
\item Per-script hardware envelopes and runtime tables for the BFS
implementations.
\end{itemize}

These items are unchanged from earlier supplementary drafts and are
recorded here as a placeholder; consult the project repository for the
detailed records.

% (the bibliography below supplies its own heading)

References to literature are as in the main paper.  The additional
references invoked by the conditional routes above are:
\begin{itemize}[leftmargin=2em,itemsep=0.2em]
\item G.~Christol, \emph{Ensembles presque p\'eriodiques
$k$-reconnaissables}, TCS 9 (1979), 141--145.
\item J.-P.~Allouche, J.~Shallit, \emph{Automatic Sequences},
Cambridge University Press, 2003.
\item B.~Adamczewski, J.~P.~Bell, \emph{On the algebraic dependence of
$p$-automatic series and $G$-functions over function fields of
positive characteristic}, J.~Number Theory 180 (2017), 559--586.
\item J.~P.~Bell, M.~Coons, \emph{A Mahler dichotomy for analytic
functions of bounded height}, J.~Number Theory 137 (2014), 91--115.
\item B.~Adamczewski, J.~P.~Bell, \emph{A problem about Mahler
functions}, Annali Sc.~Norm.~Sup.~Pisa (5) 14 (2015), 541--562.
\item B.~Adamczewski, J.~P.~Bell, M.~Coons, \emph{A height gap theorem
for coefficients of Mahler functions}, J.~Eur.~Math.~Soc.\ 2023.
\item B.~Adamczewski, C.~Faverjon, \emph{Mahler's method in several
variables and finite automata}, Compositio Math.\ 156 (2020).
\item R.~P.~Stanley, \emph{Enumerative Combinatorics}, Vol.~2,
Cambridge University Press, 1999.  Thm.~6.4.6 (algebraic $\Rightarrow$
D-finite).
\item P.~Flajolet, R.~Sedgewick, \emph{Analytic Combinatorics},
Cambridge University Press, 2009.  Thm.~VII.7.
\item Y.~Antol\'in, L.~Ciobanu, \emph{Formal conjugacy growth series
and language complexity}, Math.\ Z.\ 285 (2017), no.~3, 1015--1041.
\end{itemize}


\section{Machine verification}
\label{sec:extra-lean}

The finite and combinatorial content of the propositions above is machine-checked in Lean~4.
Two files in the Mathlib-free project carry it, so they compile in core Lean with no
dependency fetch; a clean \texttt{lake build} rechecks the whole project, including them, in
about four minutes.

\texttt{PaperExtraMod2.lean} discharges Proposition~\ref{prop:mod2-automaton} on the $39$
published terms: the closed form $u_d \bmod 2$, eventual $3$-periodicity with pre-period~$1$,
the failure of pure periodicity at $d=0$, the order-$2$ Fibonacci recurrence for $d \ge 3$,
and its failure at $d=2$. It also records the agreement with Fibonacci for $1 \le n \le 9$ and
the first deviation $u_{10}=225$. These same two recurrence facts are the mod-$2$ half of
Proposition~\ref{prop:no-recurrence-strong}(iii).

\texttt{PaperExtraCounts.lean} verifies the search grid of
Proposition~\ref{prop:no-dfinite}: that the over-determination condition
\eqref{eq:overdet} selects exactly $52$ pairs, and that every selected pair has strictly
fewer unknowns than equations. It also verifies the arithmetic behind the withdrawal in
Proposition~\ref{prop:no-recurrence-strong}(ii), namely $\binom{9}{3}=84$ monomials, $252$
unknowns and $38$ equations, and it proves the general form of the argument used in
Proposition~\ref{prop:height-gap}: for any notion of automaticity whose sequences are finitely
valued, an unbounded sequence is not automatic. The finitely-valued property appears as an
explicit hypothesis rather than being assumed silently.

\texttt{NoRecurrence.lean} certifies Proposition~\ref{prop:no-recurrence} outright. For each
order $k$ with $1 \le k \le 19$ the system of $39-k$ equations in $k$ unknowns is shown to
have no rational solution by exhibiting a Farkas witness: an integer vector $w$ with
$w^{\mathsf{T}}A = 0$ and $w^{\mathsf{T}}b \neq 0$, which forbids a solution since any
rational $c$ would give $w^{\mathsf{T}}b = w^{\mathsf{T}}(Ac) = 0$. The witnesses were
computed once in exact GMP arithmetic by the Rust tool \texttt{code/zeta\_probe/tools/norec};
Lean re-runs none of that linear algebra and verifies only the two defining integer dot
products, so a reader need not trust the search that produced them. All nineteen witnesses
are small, at most $22$ digits per entry, so the check is kernel evaluation and carries no
compiler-evaluation axiom. The Farkas step itself is proved for arbitrary integer data, so it
is visibly independent of this sequence.

One convention is worth recording, since the statement of Proposition~\ref{prop:no-dfinite}
leaves it implicit. The count $52$ is for $k \ge 1$, that is for recurrences of order at least
one. Admitting $k=0$ gives $60$ pairs and requiring $m \ge 1$ gives $50$; only
$1 \le k \le 9$, $0 \le m \le 7$ reproduces the published figure. The Lean file verifies both
counts, so the convention is fixed by the formalization rather than by the prose.

What is not formalized is stated plainly in each file. The D-finite exclusion computations of
Proposition~\ref{prop:no-dfinite}, being exact linear algebra over $43$ terms at $52$ parameter
pairs, are not re-run in Lean, and neither is the passage
from eventual periodicity to algebraicity over $\mathbb{F}_2(t)$, which needs formal power
series; nor the unboundedness of $u_d$, which is a statement about all $d$ and rests on the
growth rate of \cite{Bonfioli2}. Strict increase across the published range is verified, and
is evidence for unboundedness rather than a proof of it. The regression numbers in
Proposition~\ref{prop:height-gap} are numerics and are not formalized.

\begin{thebibliography}{9}

\bibitem{AntolinCiobanu} Y.~Antol\'in and L.~Ciobanu, \emph{Formal conjugacy growth in
acylindrically hyperbolic groups}, Int.\ Math.\ Res.\ Not.\ IMRN \textbf{2017}, no.~1,
121--157.

\bibitem{ABS} B.~Adamczewski, J.~P.~Bell and D.~Smertnig, \emph{A height gap theorem for
coefficients of Mahler functions}, J.\ Eur.\ Math.\ Soc.\ \textbf{23} (2021), 2379--2398.

\bibitem{Bonfioli1b} V.~Bonfioli, \emph{A dimension-independent collision-depth invariant for
right-corner orthoscheme reflection groups}, companion paper.

\bibitem{Christol1979} G.~Christol, \emph{Ensembles presque p\'eriodiques $k$-reconnaissables},
Theoret.\ Comput.\ Sci.\ \textbf{9} (1979), 141--145.

\bibitem{CKMR} G.~Christol, T.~Kamae, M.~Mend\`es France, G.~Rauzy,
\emph{Suites alg\'ebriques, automates et substitutions},
Bull.\ Soc.\ Math.\ France \textbf{108} (1980), 401--419.

\bibitem{AlloucheShallit} J.-P.~Allouche, J.~Shallit, \emph{Automatic Sequences: Theory,
Applications, Generalizations}, Cambridge Univ.\ Press, 2003.

\bibitem{Mahler1929} K.~Mahler, \emph{Arithmetische Eigenschaften der L\"osungen einer Klasse
von Funktionalgleichungen}, Math.\ Ann.\ \textbf{101} (1929), 342--366.

\bibitem{Nishioka} K.~Nishioka, \emph{Mahler Functions and Transcendence},
Lecture Notes in Math.\ \textbf{1631}, Springer, 1996.

\bibitem{Nesterenko1996} Yu.~V.~Nesterenko, \emph{Modular functions and transcendence
questions}, Sb.\ Math.\ \textbf{187} (1996), 1319--1348.

\bibitem{Bonfioli2} V.~Bonfioli, \emph{Transcendence of the right-triangle reflection growth
series}, companion paper; code and formalization at
\texttt{github.com/ElVec1o/pythagorean-reflection-sequence}.

\end{thebibliography}

\end{document}
